%!TEX root = ../template.tex

%% ----------------------------------------------------------------------------
%%  NOVATHESIS — 2-MainMatter/chapter-covers.tex
%% ----------------------------------------------------------------------------
%%
%% Version 7.10.7 (2026-03-25)
%% Copyright (C) 2004-26 by João M. Lourenço <joao.lourenco@fct.unl.pt>
%%
%% Reference chapter: the NOVAthesis cover page system.
%% Intended audience: developers adding new school support or refactoring
%% existing cover definitions.
%%
%% ----------------------------------------------------------------------------

\typeout{NT FILE 2-MainMatter/chapter-covers.tex}%

\chapter{The NOVAthesis Cover Page System}
\label{app:covers}

This chapter documents the cover page engine used by the \texttt{novathesis}
class.  It is written for developers who need to add cover support for a new
institution, or who want to understand or refactor an existing cover
definition.  End users who only want to configure metadata (title, author,
committee) should work exclusively with \latexinline{0-Config/3_cover.tex}.

The cover engine is implemented in:
\begin{itemize}
  \item \latexinline{NOVAthesisFiles/StyFiles/nt-covers.sty} — a modern
        \hologo{LaTeX3}/expl3 implementation that provides a high-level
        declarative \gls{API} for cover design.
\end{itemize}
The remainder of this chapter describes that \gls{API}.


%%=======================================================================
\section{Architecture Overview}
\label{sec:covers_arch}
%%=======================================================================

Every cover page is rendered through three independent layers applied in
order:

\begin{enumerate}
  \item \textbf{Background colour.}  A solid \latexinline{\pagecolor{}} fills the
        entire page.  Set with \latexinline{\thesiscover(id,bgcolor):={colorname}}.
  \item \textbf{Background image.}  A full-bleed graphic is stretched to
        \latexinline{\paperwidth} $\times$ \latexinline{\paperheight}.  Set with
        \latexinline{\thesiscover(id,image):={filename}}.
  \item \textbf{Content layer.}  A \latexinline{tikzpicture} overlay anchors a
        large \latexinline{minipage} at the top-left of the text area (offset by the
        cover margins).  Inside this minipage, elements registered with
        \latexinline{\ntaddtocover} are typeset top-to-bottom.  TikZ elements
        registered with \latexinline{tikz=true} or \latexinline{\ntAddToCoverTikz} are
        collected and rendered in a \emph{second} overlay \latexinline{tikzpicture}
        drawn after the minipage, so they escape the minipage's coordinate
        space.
\end{enumerate}

The cover text area is defined by four cover-specific margins (see
Section~\ref{sec:covers_margins}).  All \latexinline{\ntaddtocover} elements are
stacked \emph{vertically} inside this area, in the order they were
registered.


%%=======================================================================
\section{Cover Identifiers}
\label{sec:cover_ids}
%%=======================================================================

Each cover page has a two-part identifier \texttt{major\,-\,minor}:

\begin{center}
\begin{ShadedTabular}{cll}
  \toprule
  \textbf{ID} & \textbf{Description} & \textbf{When printed} \\
  % \midrule
  \texttt{1-1} & Front cover — recto (right/odd page)   & Always \\
  \texttt{1-2} & Front cover — verso (back of cover)    & When content is registered \\
  \texttt{2-1} & Front page — recto                     & When \texttt{print/frontpage=true} \\
  \texttt{2-2} & Front page — verso                     & When content is registered \\
  \texttt{N-1} & Back cover — recto (reverse of back)   & When content is registered \\
  \texttt{N-2} & Back cover — verso (outside back cover) & Always \\
  \bottomrule
\end{ShadedTabular}
\end{center}

The \emph{major} number (1, 2, N) identifies the cover pair; the \emph{minor}
number (1, 2) selects the recto or verso page within the pair.  A
comma-separated list of IDs (e.g.\ \latexinline{1-1,2-1}) can be passed to any
command that accepts cover IDs.


%%=======================================================================
\section{Setting Up the Cover Canvas}
\label{sec:covers_canvas}
%%=======================================================================

%%-----------------------------------------------------------------------
\subsection{Cover-Specific Margins}
\label{sec:covers_margins}
%%-----------------------------------------------------------------------

Margins for the text area (the minipage placed on the content layer) are
separate from the normal page margins.  They are set with the
\latexinline{\margin(cover,...):={...}} family:

\begin{latexcode}
\SetMargin*(cover,top)={8.6cm}    % default top margin for all cover pages
\SetMargin*(cover,bottom)={1.5cm}
\SetMargin*(cover,left)={3.55cm}
\SetMargin*(cover,right)={2.0cm}
\SetMargin*(cover,2-1,top)={6.6cm} % override top margin for cover 2-1 only
\end{latexcode}

The cover-ID-specific variant takes precedence when printing that particular
page.  Only the text area is affected; background images always fill the
entire paper.


%%-----------------------------------------------------------------------
\subsection{Background Colour and Image}
\label{sec:covers_bg}
%%-----------------------------------------------------------------------

These properties are stored in the \texttt{thesiscover} data family and
retrieved through a prioritised lookup that checks the current document type
and degree class before falling back to the generic key:

\begin{latexcode}
% Set a solid background colour for cover 1-1
\SetThesisCover*(1-1,bgcolor)={fctdarkblue}

% Per-doctype variant (overrides the generic key when doctype=phd)
\SetThesisCover*(phd,1-1,bgcolor)={fctdarkgreen}

% Set a background image (file must be findable by \hologo{LaTeX}/graphicx)
\SetThesisCover*(phd,1-1,image)={nova-fct-cover-green-front}

% Set the default text colour for the entire cover page
\SetThesisCover*(1-1,textcolor)={white}
\end{latexcode}

Supported keys:
\begin{itemize}
  \item \latexinline{bgcolor} — name of a colour defined with \latexinline{\definecolor}.
  \item \latexinline{image} — file name (no extension; passed to
        \latexinline{\includegraphics}).  The image is automatically scaled to
        \latexinline{\paperwidth} $\times$ \latexinline{\paperheight}.
  \item \latexinline{textcolor} — sets \latexinline{\color{...}} for the entire content
        layer of that cover page.
\end{itemize}

The lookup cascade used internally is:
\begin{latexcode}
\@DOCTYPE, \@DEGREEACR, cover-id, key
\@DOCCLASS, \@DEGREEACR, cover-id, key
\@DEGREEACR, cover-id, key
\@DEGREEACR, key
\@DOCTYPE, cover-id, key
\@DOCCLASS, cover-id, key
cover-id, key
key
\end{latexcode}

The first matching entry wins, so you can easily have a generic fallback with
targeted per-doctype overrides.


%%=======================================================================
\section{The \texttt{\textbackslash ntaddtocover} Command}
\label{sec:ntaddtocover}
%%=======================================================================

\latexinline{\ntaddtocover} is the primary interface for populating a cover page.
It registers a \emph{content block} — either a vertical box or a TikZ code
snippet — for one or more cover pages.

\medskip
\noindent\textbf{Syntax:}
\begin{latexcode}
\ntaddtocover[<attributes>]{<cover-id-list>}{<code>}
\end{latexcode}

\begin{description}
  \item[\texttt{[attributes]}] Optional comma-separated key=value list.
        See Section~\ref{sec:ntaddtocover_attrs}.
  \item[\texttt{\{cover-id-list\}}] Comma-separated list of cover IDs (e.g.\
        \latexinline{1-1,2-1}) to which this element is appended.
  \item[\texttt{\{code\}}] The content to typeset.  Evaluated inside a
        \latexinline{minipage} (unless \latexinline{tikz=true}; see
        Section~\ref{sec:ntaddtocover_tikz}).
\end{description}

Elements are rendered in the order they are registered.  The first call
produces the topmost block.

\medskip
\noindent\textbf{Minimal example — adding a title block:}
\begin{latexcode}
\ntaddtocover[halign=l, height=3cm, valign=c]{1-1}{%
  \fontsize{20}{21}\selectfont
  \textbf{\GetDocTitle(\@LANG@COVER,main)}%
}
\end{latexcode}


%%-----------------------------------------------------------------------
\subsection{Attribute Reference}
\label{sec:ntaddtocover_attrs}
%%-----------------------------------------------------------------------

All attributes are optional.  Defaults are shown in parentheses.

\medskip
\noindent\textbf{Filter attributes} (the element is silently skipped when
the filter does not match the current document):

\begin{description}
  \item[\texttt{status=\{...\}}] Comma-separated list of document statuses
        for which this element should appear.  Valid values: \texttt{working},
        \texttt{provisional}, \texttt{final}.  Empty (default) means
        \emph{always}.
  \item[\texttt{degree=\{...\}}] Comma-separated list of document types
        for which this element should appear.  Valid values: \texttt{phd},
        \texttt{msc}, \texttt{bsc}, \texttt{plain}, \texttt{phdplan},
        \texttt{phdprop}, \texttt{mscplan}.  Empty (default) means
        \emph{always}.
\end{description}

\medskip
\noindent\textbf{Box geometry attributes:}

\begin{description}
  \item[\texttt{width=\{<dim>\}}] Width of the box (default: full text area
        width minus \texttt{hspace}).  Any valid \TeX\ dimension.
  \item[\texttt{height=\{<dim>\}}] Fixed height of the box.  When omitted, the
        box height matches its natural content height.  Using a fixed height
        together with \texttt{valign} lets you vertically place content inside
        a reserved area.
  \item[\texttt{height/phd=\{<dim>\}}] Same as \texttt{height}, but only
        applied when the document type is \gls{phd}.  Overrides \texttt{height} for
        \gls{phd} documents.
  \item[\texttt{height/msc=\{<dim>\}}] Same, but for \gls{msc} documents.
  \item[\texttt{minheight=\{<dim>\}}] Minimum height of the box, even when the
        content is shorter (default: \texttt{0pt}).
  \item[\texttt{hspace=\{<dim>\}}] Horizontal offset from the left margin
        (default: \texttt{0pt}).  Useful for indenting elements that must not
        start at the margin.
\end{description}

\medskip
\noindent\textbf{Alignment attributes:}

\begin{description}
  \item[\texttt{halign=c|l|r|j}] Horizontal alignment of content inside the
        box (default: \texttt{c}).  Values: \texttt{c} = \latexinline{\centering},
        \texttt{l} = \latexinline{\raggedright}, \texttt{r} = \latexinline{\raggedleft},
        \texttt{j} = justified (no alignment command).
  \item[\texttt{valign=c|t|b}] Vertical alignment of content inside a
        fixed-height box (default: \texttt{c}).  Ignored when no \texttt{height}
        is given.
\end{description}

\medskip
\noindent\textbf{Colour attributes:}

\begin{description}
  \item[\texttt{color=\{<colorname>\}}] Foreground text colour inside the box.
        Overrides \texttt{textcolor} from \latexinline{\thesiscover}.
  \item[\texttt{bgcolor=\{<colorname>\}}] Background colour of the box
        (applied via \latexinline{\colorbox}).  Default: transparent.
\end{description}

\medskip
\noindent\textbf{Vertical spacing attributes:}

The \texttt{vspace} attribute controls the vertical gap \emph{before} the
element.  It accepts either a \TeX\ dimension or a plain integer:

\begin{description}
  \item[\texttt{vspace=\{<dim>\}}] Insert a fixed \latexinline{\vspace*{<dim>}} above
        the element.
  \item[\texttt{vspace=\{<n>\}}] When the value is a plain integer, insert an
        elastic \latexinline{\vskip 0pt plus <n>fill} instead.  This creates a
        stretchable gap — setting \texttt{vspace=1} on consecutive elements
        distributes the available vertical space evenly between them.
  \item[\texttt{vspace/phd}, \texttt{vspace/msc}] Per-doctype override of
        \texttt{vspace}.  Takes precedence over the generic \texttt{vspace}
        when the document type matches.
  \item[\texttt{vspace/1-1}, \texttt{vspace/2-1}] Per-cover-ID override.
        Takes precedence when the element is being rendered on that specific
        cover.
  \item[\texttt{vspace/phd/1-1}, \texttt{vspace/msc/1-1}, etc.] Combined
        per-doctype per-cover override.  Highest precedence.
\end{description}

\medskip
\noindent\textbf{Layout control attributes:}

\begin{description}
  \item[\texttt{newpar=true|false}] Emit a \latexinline{\par} after the element
        (default: \texttt{true}).  Set \texttt{newpar=false} to place the
        next element on the same line.
  \item[\texttt{tikz=true|false}] Treat the content block as raw TikZ code to
        be injected into the overlay TikZ picture rather than typeset inside
        the minipage (default: \texttt{false}).  See
        Section~\ref{sec:ntaddtocover_tikz}.
\end{description}


%%-----------------------------------------------------------------------
\subsection{Stacking Elements Vertically}
\label{sec:covers_vstack}
%%-----------------------------------------------------------------------

The most common cover layout is a vertical stack of boxes, similar to a
\latexinline{minipage} column.  Each \latexinline{\ntaddtocover} call appends one box to the
column.  The following design rules apply:

\begin{enumerate}
  \item The element appears immediately after the previous one (zero gap by
        default).
  \item Use \texttt{vspace=\{<dim>\}} for a fixed gap and
        \texttt{vspace=\{<n>\}} (integer) for a stretchable gap.
  \item A single \texttt{vspace=1} between two groups of elements pushes one
        group to the top and the other to the bottom.  Use equal integer
        weights (e.g.\ two elements with \texttt{vspace=1}) for equal
        distribution.
  \item Use \texttt{height} to reserve vertical space even when content may
        vary in size (e.g.\ optional subtitle lines).  Combine with
        \texttt{valign=c} to centre the content vertically inside the reserved
        area.
\end{enumerate}

Example — logo, title, elastic gap, date:
\begin{latexcode}
% Logo at the top of the text area
\ntaddtocover[halign=l]{1-1,2-1}{%
  \includegraphics[height=1.5cm]{my-school-logo}%
}

% Title — reserved box of fixed height so subtitle does not shift the author
\ntaddtocover[vspace=5mm, halign=l, height=4cm, valign=c]{1-1,2-1}{%
  \fontsize{18}{20}\selectfont
  \textbf{\GetDocTitle(\@LANG@COVER,main)}%
}

% Author
\ntaddtocover[halign=l]{1-1,2-1}{%
  \fontsize{14}{14}\selectfont
  \GetDocAuthor(name)%
}

% Elastic gap pushes the date to the bottom of the text area
\ntaddtocover[vspace=1, halign=l]{1-1,2-1}{%
  \fontsize{11}{11}\selectfont
  \PrintDateISO{\GetDocDate(submission,year)-\GetDocDate(submission,month)}%
}
\end{latexcode}

\medskip

\noindent\textbf{Conditional content per document type:}

Use the \texttt{degree} filter or the internal \latexinline{\ifphddoc}/\latexinline{\ifmscdoc}
helpers when the cover layout differs between \gls{phd} and \gls{msc} documents:

\begin{latexcode}
\ntaddtocover[halign=l, height=3cm, valign=c]{1-1}{%
  \ifphddoc{%
    \fontsize{20}{21}\selectfont
    \textbf{\GetDocTitle(\@LANG@COVER,main)}%
  }{%
    \fontsize{17}{17}\selectfont
    \textbf{\GetDocAuthor(name)}%
  }%
}
\end{latexcode}

Alternatively, register two separate elements with complementary
\texttt{degree} filters and identical heights so they occupy the same space:

\begin{latexcode}
\ntaddtocover[degree=phd, height=3cm, valign=c]{1-1}{%
  \fontsize{20}{21}\selectfont \textbf{\GetDocTitle(\@LANG@COVER,main)}%
}
\ntaddtocover[degree=msc, height=3cm, valign=c]{1-1}{%
  \fontsize{17}{17}\selectfont \textbf{\GetDocAuthor(name)}%
}
\end{latexcode}


%%-----------------------------------------------------------------------
\subsection{Embedding TikZ Code via \texttt{tikz=true}}
\label{sec:ntaddtocover_tikz}
%%-----------------------------------------------------------------------

Setting \texttt{tikz=true} in the attribute list causes the content block to
be collected as \emph{overlay TikZ code} rather than box content.  The block
is executed inside a \latexinline{tikzpicture[remember picture, overlay]} that covers
the whole page, drawn after the text-layer minipage has been placed.

This means the coordinate origin is the \latexinline{current page} node, not the
text area, so you can position elements anywhere on the page with absolute
coordinates.  The \latexinline{calc} TikZ library (\latexinline{\usetikzlibrary{calc}}) is
loaded by default and can be used freely.

\begin{latexcode}
% Place the university logo at the top-left corner (absolute coordinates)
\ntaddtocover[tikz]{1-1}{%
  \node[anchor=center]
    at ($(current page.north west) + (20mm,-20mm)$)
    {\includegraphics[width=21mm]{\GetUniversity(logo,neg)}};%
  % Vertical rule separator
  \draw[thick]
    ($(current page.north west) + (3.5cm,-1cm)$) -- ++(0,-2cm);%
}
\end{latexcode}

All \texttt{tikz=true} blocks registered for the same cover are concatenated
into a single \latexinline{tikzpicture}, in registration order.

\medskip
\noindent\textbf{When to use \texttt{tikz=true} vs.\ plain box elements:}
\begin{itemize}
  \item Use plain box elements for text content whose position depends on the
        height of other elements (titles, author names, committee lists).
        The vertical-stacking model handles reflow automatically.
  \item Use \texttt{tikz=true} for graphics or decorations that need to be
        anchored at a fixed absolute position (logos at a corner, background
        watermarks, ruled lines).
\end{itemize}


%%=======================================================================
\section{The \texttt{\textbackslash ntAddToCoverTikz} Command}
\label{sec:ntaddtocovertikz}
%%=======================================================================

\latexinline{\ntAddToCoverTikz} is a convenience command for registering individual
TikZ nodes or drawing commands.  It mirrors the syntax of TikZ's own node
and draw commands, inserting them into the same overlay picture as
\latexinline{\ntaddtocover[tikz]}.

\medskip
\noindent\textbf{Syntax:}
\begin{latexcode}
\ntAddToCoverTikz<cmd>[opts](position)[rest]{cover-id-list}{code}
\end{latexcode}

\begin{description}
  \item[\texttt{<cmd>}] Optional: a TikZ command other than \latexinline{\node}
        (use angle brackets, e.g.\ \latexinline{<\draw>}).  When omitted, the
        command defaults to \latexinline{\node}.
  \item[\texttt{[opts]}] Optional TikZ options for the node or command (square
        brackets, default empty).
  \item[\texttt{(position)}] Required: position of the node or command, in
        the TikZ coordinate system relative to \latexinline{current page} (parentheses).
  \item[\texttt{[rest]}] Optional: additional TikZ node syntax placed between
        the position and the final \latexinline{{code}} argument (square brackets,
        default empty).
  \item[\texttt{\{cover-id-list\}}] Comma-separated list of cover IDs.
  \item[\texttt{\{code\}}] Node content (when using the default \latexinline{\node}
        command) or the draw-path code (when using a custom command).
\end{description}

\medskip
\noindent\textbf{How the arguments map to TikZ output:}

When no command is specified, the generated TikZ is:
\begin{latexcode}
\node[opts] at (position) rest {code};
\end{latexcode}

When a command \latexinline{<\draw>} is specified:
\begin{latexcode}
\draw[opts] (position) rest;
\end{latexcode}
(note: \texttt{code} is ignored in command mode)

\medskip
\noindent\textbf{Examples:}
\begin{latexcode}
% Default \node — place a logo
\ntAddToCoverTikz
  [anchor=north east]
  (current page.north east)
  {1-1,2-1}
  {\includegraphics[width=4cm]{my-logo}}

% Custom command: \draw
\ntAddToCoverTikz<\draw>
  [thick, white]
  ($(current page.north west) + (1cm,-4cm)$)
  [-- ++(\paperwidth-2cm, 0)]
  {1-1,2-1}
  {}
\end{latexcode}

\medskip
\noindent\textbf{Note on \texttt{yscale=-1}.}  The overlay picture for
\latexinline{\ntAddToCoverTikz} uses \latexinline{yscale=-1}, so the $y$-axis points
\emph{downward}.  Coordinates like \latexinline{(current page.north west) + (x,-y)}
move right and \emph{down} from the top-left corner.  The plain
\latexinline{\ntaddtocover[tikz]} blocks use a \emph{standard} overlay without
\texttt{yscale}, so their coordinate conventions differ.


%%=======================================================================
\section{Cover Hooks}
\label{sec:covers_hooks}
%%=======================================================================

Hooks let you execute code at specific points in the cover rendering pipeline
without modifying the \latexinline{\ntaddtocover} element list.  They are registered
with \latexinline{\NTAddToHook}:

\begin{latexcode}
\NTAddToHook{cover/1-1/pre}{\sffamily\color{white}}
\NTAddToHook{cover/2-1/pre}{\sffamily}
\NTAddToHook{cover/post}{\pagecolor{white}} % reset after each cover
\end{latexcode}

Available hook points (in execution order for a single cover page):

\begin{center}
\begin{ShadedTabular}{ll}
  \toprule
  \textbf{Hook name}              & \textbf{When executed} \\
  % \midrule
  \texttt{cover/pre}              & Before any cover page \\
  \texttt{cover/<id>/pre}         & Before cover \texttt{<id>} \\
  \texttt{cover/<id>/text/pre}    & Before text-layer elements \\
  \texttt{cover/<id>/text/post}   & After text-layer elements \\
  \texttt{cover/<id>/bgcolor/pre} & Before background colour is applied \\
  \texttt{cover/<id>/bgcolor/post}& After background colour is applied \\
  \texttt{cover/<id>/image/pre}   & Before background image is rendered \\
  \texttt{cover/<id>/image/post}  & After background image is rendered \\
  \texttt{cover/<id>/grid/pre}    & Before debug grid (only in debug mode) \\
  \texttt{cover/<id>/grid/post}   & After debug grid \\
  \texttt{cover/<id>/post}        & After cover \texttt{<id>} \\
  \texttt{cover/post}             & After any cover page \\
  \bottomrule
\end{ShadedTabular}
\end{center}

\medskip
\noindent\textbf{Typical use cases:}
\begin{itemize}
  \item Switch to a sans-serif font or set a global colour for the whole cover
        page inside the \texttt{cover/<id>/pre} hook.
  \item Reset font or colour in \texttt{cover/post} to avoid bleeding into the
        rest of the document.
  \item Load cover-specific fonts (e.g.\ via \latexinline{\setmainfont} for
        \hologo{LuaLaTeX}/\hologo{XeLaTeX}) in \texttt{cover/pre}.
\end{itemize}


%%=======================================================================
\section{Debugging a Cover Layout}
\label{sec:covers_debug}
%%=======================================================================

Enable debug mode in \latexinline{0-Config/1_novathesis.tex}:

\begin{latexcode}
\ntsetup{debug={cover,spine}}
\end{latexcode}

With \texttt{cover} in the debug set, the engine automatically calls
\latexinline{\debuggrid} over every cover page and draws thin frames around each
text box placed by \latexinline{\ntaddtocover}.

You can also call \latexinline{\debuggrid} manually inside a cover element:
\begin{latexcode}
\ntaddtocover{1-1}{\debuggrid[cyan]}
\end{latexcode}

The grid draws:
\begin{itemize}
  \item 1\,mm ticks (lightest),
  \item 5\,mm ticks,
  \item 1\,cm ticks,
  \item 5\,cm ticks (darkest).
\end{itemize}

All lines respect the cover margins, so the grid origin is the top-left
corner of the text area.

\medskip
A fast iteration workflow:
\begin{bashcode}
# Build the cover only (mode=2 skips all document content)
python .Build/build.py nova/fct --mode 2
\end{bashcode}

This rebuilds only the front cover page, giving fast feedback when tuning
positions and margins.


%%=======================================================================
\section{Cover Rendering: Execution Order Summary}
\label{sec:covers_order}
%%=======================================================================

The following sequence describes what happens when
\latexinline{\ntprintcoversingle{1}{1}} (cover \texttt{1-1}) is called:

\begin{enumerate}
  \item A fresh group is opened; \texttt{cover/pre} and
        \texttt{cover/1-1/pre} hooks are run.
  \item A \gls{PDF} bookmark is inserted.
  \item \texttt{thispagestyle{empty}} suppresses headers and footers.
  \item Background colour lookup: the engine searches the
        \texttt{thesiscover} data store for a matching \texttt{bgcolor} key
        and calls \latexinline{\pagecolor}.
  \item Background image lookup: if an \texttt{image} key is found, a
        full-page \latexinline{tikzpicture} overlay renders it.
  \item \texttt{cover/1-1/text/pre} hook.
  \item Cover margins are read; a \latexinline{tikzpicture} overlay anchors a
        \latexinline{minipage} at \latexinline{current page.north west} offset by those
        margins.
  \item Inside the minipage, each element registered with
        \latexinline{\ntaddtocover[…]{1-1}} is processed in order.
        For each element the engine: resolves \texttt{vspace}, typeset the
        content into a box, applies \texttt{bgcolor} wrapping, and emits the
        box.
  \item \texttt{tikz=true} elements accumulated during step~8 are emitted in
        a second overlay \latexinline{tikzpicture}.
  \item \texttt{cover/1-1/text/post}, \texttt{cover/1-1/post}, and
        \texttt{cover/post} hooks are run.
  \item The group is closed.
\end{enumerate}


%%=======================================================================
\section{Step-by-Step: Adding Cover Support for a New School}
\label{sec:covers_new_school}
%%=======================================================================

This section walks through adding a complete cover definition.  Substitute
\texttt{myuniv} and \texttt{myfac} with the institution identifiers.

\medskip
\noindent\textbf{1.  Create the directory tree:}
\begin{bashcode}
NOVAthesisFiles/Schools/myuniv/myfac/
NOVAthesisFiles/Schools/myuniv/myfac/Images/
\end{bashcode}

\medskip
\noindent\textbf{2.  Create the defaults file}
\latexinline{myuniv-myfac-defaults.clo}:

\begin{latexcode}
%!TEX root=../../../template.tex
%% myuniv/myfac/myuniv-myfac-defaults.clo

\typeout{NT FILE schools/myuniv/myfac/myuniv-myfac-defaults.clo}%


%%=======================================================
%% COLOURS
%%=======================================================

\definecolor{myfacblue}{RGB}{0, 64, 128}
\definecolor{myfacgold}{RGB}{190, 155, 80}


%%=======================================================
%% INSTITUTION STRINGS  (one line per supported language)
%%=======================================================

\SetUniversity*(en)={My University}
\SetUniversity*(pt)={Minha Universidade}
\SetSchool*(en)={My Faculty}
\SetSchool*(pt)={Minha Faculdade}

\SetSchool*(logo,en)={myuniv-myfac-logo-en}
\SetSchool*(logo,pt)={myuniv-myfac-logo-pt}
\IfMemVoidTF{School}{logo,\@LANG@COVER}
  {\SetSchool*(logo)={\GetSchool(logo,en)}}
  {\SetSchool*(logo)={\GetSchool(logo,\@LANG@COVER)}}


%%=======================================================
%% COMMITTEE ORDER
%%=======================================================

\SetCommitteeOrder*()={c,r,a,ca,m}%


%%=======================================================
%% LAYOUT
%%=======================================================

\SetMargin*(cover,top)={4.0cm}
\SetMargin*(cover,bottom)={2.0cm}
\SetMargin*(cover,left)={2.5cm}
\SetMargin*(cover,right)={2.5cm}


%%=======================================================
%% COVER BACKGROUNDS
%%=======================================================

\SetThesisCover*(1-1,bgcolor)={myfacblue}
\SetThesisCover*(1-1,textcolor)={white}
\SetThesisCover*(2-1,bgcolor)={white}
\SetThesisCover*(N-2,bgcolor)={myfacblue}


%%=======================================================
%% COVER ELEMENTS — cover 1-1 (front cover)
%%=======================================================

% Logo — absolutely positioned in the top-right area
\ntaddtocover[tikz]{1-1}{%
  \node[anchor=north east]
    at ($(current page.north east) + (-1cm,-1cm)$)
    {\includegraphics[height=2cm]{\GetSchool(logo)}};%
}

% Title
\ntaddtocover[halign=l, height=4cm, valign=c]{1-1}{%
  \fontsize{20}{22}\selectfont\bfseries
  \GetDocTitle(\@LANG@COVER,main)%
}

% Author
\ntaddtocover[halign=l]{1-1}{%
  \fontsize{14}{14}\selectfont
  \GetDocAuthor(name)%
}

% Degree name + date pushed to the bottom
\ntaddtocover[vspace=1, halign=l]{1-1}{%
  \fontsize{11}{11}\selectfont
  \GetDegreeName(\@LANG@COVER)\\
  \PrintDateISO{\GetDocDate(submission,year)-%
                \GetDocDate(submission,month)}%
}


%%=======================================================
%% COVER ELEMENTS — cover 2-1 (front page)
%%=======================================================

% Reproduce key elements on the front page
\ntaddtocover[halign=l]{2-1}{%
  \includegraphics[height=1.5cm]{\GetSchool(logo)}\\[4mm]%
  \fontsize{16}{18}\selectfont\bfseries
  \GetDocTitle(\@LANG@COVER,main)%
}

% Advisers
\ntaddtocover[vspace=2, halign=l]{2-1}{%
  \MemSet{adviserstringfont}{}{\bfseries}%
  \fontsize{11}{11}\selectfont
  \ntprintpersons{0.9}{2}{adviser}{\GetAdviserOrder()}%
}

% Committee (final status only)
\ntaddtocover[vspace=1, halign=l, status={final}]{2-1}{%
  \MemSet{committeetitlestringfont}{}{\bfseries}%
  \fontsize{10}{10}\selectfont
  \ntprintpersons{0.9}{2}{committee}{\GetCommitteeOrder()}%
}

% Date at the bottom
\ntaddtocover[vspace=1, halign=l]{2-1}{%
  \fontsize{11}{11}\selectfont
  \GetDegreeName(\@LANG@COVER)\\
  \PrintDateISO{\GetDocDate(submission,year)-%
                \GetDocDate(submission,month)}%
}
\end{latexcode}

\medskip
\noindent\textbf{3.  Add logos} to
\latexinline{NOVAthesisFiles/Schools/myuniv/myfac/Images/}.
Supported formats: \texttt{.pdf} (preferred for vector logos), \texttt{.png},
\texttt{.jpg}.

\medskip
\noindent\textbf{4.  Test with cover-only mode:}
\begin{bashcode}
python .Build/build.py myuniv/myfac --mode 2
\end{bashcode}

\medskip
\noindent\textbf{5.  Test all document types:}
\begin{bashcode}
python .Build/build.py myuniv/myfac --doctype phd --docstatus final
python .Build/build.py myuniv/myfac --doctype msc --docstatus final
\end{bashcode}

\medskip
\noindent\textbf{6.  Enable debug mode} during development to verify spacing:
\begin{latexcode}
\ntsetup{debug={cover}}
\end{latexcode}


%%=======================================================================
\section{Real-World Examples}
\label{sec:covers_examples}
%%=======================================================================

%%-----------------------------------------------------------------------
\subsection{NOVA/FCT: Mixed \hologo{LaTeX} and TikZ}
\label{sec:covers_nova_fct}
%%-----------------------------------------------------------------------

The \gls{FCT} cover (\latexinline{NOVAthesisFiles/Schools/nova/fct/nova-fct-defaults.clo})
demonstrates a typical mixed strategy:

\begin{itemize}
  \item All logos and graphical decorations (university logo, school logo,
        separator rule, department name) are placed with
        \latexinline{\ntaddtocover[tikz]{...}}, giving pixel-precise positioning
        independent of the text flow.
  \item All text content (title, author, degree, date) uses plain
        \latexinline{\ntaddtocover} elements stacked in the cover's text area, which
        starts at \texttt{margin(cover,top)=8.6cm} — well below the
        absolutely-positioned logo band.
  \item The front page (cover \texttt{2-1}) reuses the same elements but with
        a lighter margin (\texttt{6.6cm} top) and adds advisers and committee.
  \item \gls{phd} and \gls{msc} documents swap the order of title and author name using
        \latexinline{\ifphddoc} conditionals inside shared elements.
\end{itemize}

Key cover-margin setting:
\begin{latexcode}
\SetMargin*(cover,top)={8.6cm}   % leaves room for the TikZ logo band
\SetMargin*(cover,2-1,top)={6.6cm} % front page has a smaller logo band
\end{latexcode}

Background setup (colour varies by document class):
\begin{latexcode}
\def\@novafctcovercolor{red}          % default (BSc/other)
\ifdocclass{msc}{\def\@novafctcovercolor{blue}}{}
\ifdocclass{phd}{\def\@novafctcovercolor{green}}{}
\SetThesisCover*(\@DOCTYPE,1-1,image)=
    {nova-fct-cover-\@novafctcovercolor-front}
\end{latexcode}


%%-----------------------------------------------------------------------
\subsection{ULISBOA/IST: Pure \hologo{LaTeX} Box Stacking}
\label{sec:covers_ist}
%%-----------------------------------------------------------------------

The IST cover
(\latexinline{NOVAthesisFiles/Schools/ulisboa/ist/ulisboa-ist-defaults.clo})
uses no TikZ for content elements at all.  Every element — logo, school
name, a decorative image, title, author, advisers, committee, date — is
placed as a plain box element with explicit \texttt{height} and
\texttt{vspace} values.  This approach is simpler to reason about and
produces a layout that tolerates variable-length content gracefully.

The layout sequence for cover \texttt{1-1}:
\begin{latexcode}
\ntaddtocover[halign=l, hspace=-2.25mm]{1-1,2-1}{logo image}
\ntaddtocover[vspace=5mm, hspace=23mm]{1-1,2-1}{school name text}
\ntaddtocover[vspace=1.62cm, height=5cm]{1-1,2-1}{decorative image}
\ntaddtocover[vspace=1.2cm,  height=1.35cm, valign=c]{1-1,2-1}{title}
\ntaddtocover[vspace=1.0cm,  height=1.05cm]{1-1,2-1}{author name}
\ntaddtocover[vspace=1.7cm,  height=3.0cm, valign=t]{1-1}{advisers}
\ntaddtocover[vspace=1,      height=1.6cm, valign=t]{1-1}{degree text}
\ntaddtocover[vspace=1, height=1.25cm, valign=t,
              status={working,provisional}]{1-1}{draft message}
\ntaddtocover[vspace=1, height=1.25cm, valign=t,
              status={final}]{1-1}{approved message}
\ntaddtocover[vspace=1, height=8.5mm, valign=b]{1-1,2-1}{date}
\end{latexcode}

Each \texttt{height} reservation ensures that regardless of whether the
optional subtitle, specialization line, or committee is present, the
elements below remain at a predictable vertical position.


%%=======================================================================
\section{Quick Reference}
\label{sec:covers_quickref}
%%=======================================================================

\begin{center}
\renewcommand{\arraystretch}{1.3}
\begin{ShadedTabular}{lp{8cm}}
  \toprule
  \textbf{Command / Setting} & \textbf{Purpose} \\
  % \midrule
  \latexinline{\margin(cover,top):={...}}          & Cover text-area top margin \\
  \latexinline{\margin(cover,<id>,top):={...}}     & Per-cover override \\
  \latexinline{\thesiscover(<id>,bgcolor):={...}}  & Page background colour \\
  \latexinline{\thesiscover(<id>,image):={...}}    & Full-page background image \\
  \latexinline{\thesiscover(<id>,textcolor):={...}}& Default text colour \\
  \latexinline{\ntaddtocover[attrs]{ids}{code}}    & Register a text box element \\
  \latexinline{\ntaddtocover[tikz]{ids}{code}}     & Register TikZ overlay code \\
  \latexinline{\ntAddToCoverTikz<cmd>[opts](…){ids}{code}} & Register a TikZ node/command \\
  \latexinline{\NTAddToHook{cover/<id>/pre}{...}}  & Inject code before a cover \\
  \latexinline{\debuggrid[color]}                  & Render a measurement grid \\
  \latexinline{\ntsetup{debug={cover}}}            & Enable cover debug mode \\
  \latexinline{python .Build/build.py ... --mode 2} & Build cover pages only \\
  \bottomrule
\end{ShadedTabular}
\end{center}

